\section{Galois connections}

\subsection{Dominio concreto}

Il dominio concreto $C$ rappresenta l'universo dei valori reali che il
programma può assumere.
In questo caso viene considerato il dominio concreto come l'insieme delle parti
dell'insieme dei numeri interi.

\begin{definizione}{Lattice del Dominio Concreto}{concrete_domain_lattice}
Il dominio concreto $\langle \wp(\mathbb{Z}), \subseteq, \cap, \cup, \mathbb{Z}, \emptyset \rangle$
forma un reticolo completo (\Cref{def:complete_lattice}), dove:
\begin{itemize}
    \item L'ordine parziale è l'inclusione insiemistica ($\subseteq$).
    \item $S_1 \subset S_2$ indica che $S_1$ è ``più preciso'' di $S_2$.
    \item $S_1 \nsubseteq S_2 \land S_2 \nsubseteq S_1$ indica che $S_1$ e
    $S_2$ non sono comparabili.
    \item Il join ($\cup$) rappresenta l'unione (o ``incertezza''), $x$ deve
    essere in $S_1$ oppure in $S_2$.
    \item Il meet ($\cap$) rappresenta l'intersezione, $x$ deve essere in
    $S_1$ e in $S_2$.
    \item Il Top ($\top = \mathbb{Z}$) rappresenta ``qualsiasi valore''
    (massima incertezza, es. gli input esterni).
    \item Il Bottom ($\bot = \emptyset$) rappresenta ``nessun valore''
    (codice irraggiungibile).
\end{itemize}

\begin{figure}[H]
    \centering
    \includegraphics[width=0.4\textwidth]{images/th_07/01.png}
    \caption{L'immagine illustra il reticolo completo del dominio concreto.}
    \label{fig:th_07_01}
\end{figure}

\end{definizione}

\subsubsection{Dominio astratto}
Il dominio astratto $A$ traccia un tipo specifico di informazione
(es. segni, intervalli, parità).
Esiste una relazione formale tra un dominio concreto $C$ e un dominio astratto
$A$ mediata da due funzioni:
\begin{itemize}
    \item Funzione di astrazione $\alpha: C \rightarrow A$: mappa un insieme
    di valori concreti nel corrispondente valore astratto.
    \item Funzione di concretizzazione $\gamma: A \rightarrow C$: mappa un
    valore astratto nell'insieme di valori concreti che esso rappresenta.
\end{itemize}

\begin{nota}{Gerarchie tra domini}{}
I domini astratti possono essere organizzati in una gerarchia e possono essere
combinati per formare domini più complessi.
Questi argomenti verranno trattati in seguito.
\end{nota}

\subsection{Intuizione sulla soundness dei domini astratti}

Affinché un dominio astratto sia corretto (sound), il passaggio tra
concreto e astratto deve preservare la correttezza dell'informazione.

Partendo da un valore concreto $c \in C$:
\begin{enumerate}
    \item L'astrazione: $a = \alpha(c)$.
    \item La concretizzazione: $c' = \gamma(a) = \gamma(\alpha(c))$.
\end{enumerate}
Per avere soundness, il risultato $c'$ deve essere una sovra-approssimazione
(o uguale) del valore di partenza $c$.
Ovvero, nel processo di astrazione si può perdere precisione (ottenere
un insieme più grande), ma non si possono perdere valori originali.

Analogamente, partendo da un valore astratto $a \in A$:
\begin{enumerate}
    \item La concretizzazione: $c = \gamma(a)$.
    \item La ri-astrazione: $a' = \alpha(c) = \alpha(\gamma(a))$.
\end{enumerate}
Il risultato $a'$ deve essere più preciso (o uguale) rispetto ad
$a$ ($a' \sqsubseteq_A a$). Concretizzare non perde informazione, quindi
ri-astrarre dovrebbe restituire qualcosa di idealmente identico o migliore.

\subsection{Monotonia}

Le funzioni $\alpha$ e $\gamma$ devono essere monotone (devono preservare
l'ordine).

\begin{teorema}{Monotonia di $\alpha$}{alpha_monotonicity}
$\alpha$ è monotona: $c_1 \subseteq c_2 \Rightarrow \alpha(c_1) \sqsubseteq_A \alpha(c_2)$.

Se si astrae qualcosa che rappresenta meno informazione (insieme più piccolo),
si ottiene un elemento astratto con meno o uguale informazione.
\end{teorema}

\begin{teorema}{Monotonia di $\gamma$}{gamma_monotonicity}
$\gamma$ è monotona: $a_1 \sqsubseteq_A a_2 \Rightarrow \gamma(a_1) \subseteq \gamma(a_2)$.

Se si concretizza qualcosa di più preciso, si ottiene un insieme concreto
più piccolo.
\end{teorema}

\subsection{Galois Connection}

La Galois Connection (connessione di Galois) è la formalizzazione algebrica
che lega dominio concreto e astratto garantendo le proprietà di soundness e
best abstraction.

\begin{definizione}{Galois Connection}{galois_connection}
Siano $\langle C, \sqsubseteq_C \rangle$ e $\langle A, \sqsubseteq_A \rangle$
due posets. Due funzioni monotone $\alpha: C \rightarrow A$ e
$\gamma: A \rightarrow C$ formano una Galois Connection se soddisfano le
seguienti proprietà:
\begin{itemize}
\item $\gamma \circ \alpha$ è estensiva: $\forall c \in C . c \subseteq_C \gamma(\alpha(c))$.
\item $\alpha \circ \gamma$ è riduttiva: $\forall a \in A . \alpha(\gamma(a)) \sqsubseteq_A a$.
\end{itemize}
\end{definizione}

\begin{nota}{Interpretazione delle proprietà di GC}{gc_interpretation}
Le proprietà di estensività e riduttività garantiscono che:
\begin{itemize}
    \item Quando si astrae e poi concretizza ($\gamma \circ \alpha$), si
    ottiene una sovra-approssimazione del valore originale (perdita di
    precisione).
    \item Quando si concretizza e poi si astrae, si ottiene un valore astratto
    più preciso o uguale (nessuna perdita di informazione).
\end{itemize}
\end{nota}


\subsection{Galois insertion (GI)}

Una Galois Connection (\Cref{def:galois_connection}) può contenere elementi
nell'astratto che sono ridondanti (più elementi astratti rappresentano lo
stesso oggetto concreto)
Quando viene rimossa questa ridondanza, si ottiene una Galois Insertion
(Intersezione di Galois).

\begin{definizione}{Galois Insertion}{galois_insertion}

    Siano $\langle C, \sqsubseteq_C \rangle$ e $\langle A, \sqsubseteq_A \rangle$
due posets e $\alpha: C \rightarrow A$ e $\gamma: A \rightarrow C$ due funzioni
monotone. Supponendo che esse formino una Galois Connection
(\Cref{def:galois_connection}) tra $C$ e $A$, se $\alpha \circ \gamma$ è
l'identità su $A$ questa si chiama Galois Insertion.

\end{definizione}

\begin{nota}{Riduzione da GC a GI}{gc_to_gi}
Ogni GC può essere ridotta a una GI rimuovendo gli elementi ridondanti
dal dominio astratto (domain reduction).
Se non si ha una GI, significa che esistono almeno due elementi astratti che
rappresentano la stessa informazione concreta.
\end{nota}

\begin{figure}[H]
    \centering
    \begin{tikzpicture}
        \node[draw] (C) at (0,0) {Concreto};
        \node[draw] (A) at (4,0) {Astratto};
        \draw[->, bend left] (C) to node[above] {$\alpha$} (A);
        \draw[->, bend left] (A) to node[below] {$\gamma$} (C);
        \node[below=1cm of C, align=center, font=\small] {GC generica:\\$\alpha(\gamma(a)) \sqsubseteq a$};
        \node[below=1cm of A, align=center, font=\small] {GI:\\$\alpha(\gamma(a)) = a$};
    \end{tikzpicture}
    \caption{Relazione tra domini in una GC e in una GI.}
    \label{fig:gc_vs_gi}
\end{figure}

\subsection{Best abstraction}

Siano $\langle C, \subseteq_C \rangle$ e $\langle A, \sqsubseteq_A \rangle$
due posets, e $\alpha: C \rightarrow A$ e $\gamma: A \rightarrow C$ due
funzioni monotone. Queste formano una Connessione di Galois, se:
\[
\forall c \in C, \forall a \in A . \alpha(c) \sqsubseteq_A a \iff c \sqsubseteq_C \gamma(a)
\]
dove:
\begin{itemize}
    \item $\alpha(c) \sqsubseteq_A a \implies c \sqsubseteq_C \gamma(a)$ assicura
    la soundess
    \item $c \sqsubseteq_C \gamma(a) \implies \alpha(c) \sqsubseteq_A a$ assicura
    la best abstraction.
\end{itemize}

\subsubsection{Problema della best abstraction}
Non sempre esiste una Galois Connection tra un dominio concreto e un
dominio astratto arbitrario.

\begin{esempio}{Poliedri convessi e cerchi}{convex_polyhedra_circle}
Considerando l'insieme $P = \{(x,y) | x^2 + y^2 \le 1\}$ (un cerchio).
Se il dominio astratto è quello dei poliedri convessi ($a \cdot x + b \cdot y \le c$),
non esiste la best abstraction per il cerchio.

È approssimare il cerchio con un quadrato, un ottagono, un dodecagono, etc.,
ma è sempre possibile trovare un poliedro con più lati che approssima il
cerchio meglio del precedente senza mai coincidere con il cerchio stesso.

\begin{figure}[H]
    \centering
    \includegraphics[width=0.3\textwidth]{images/th_07/02.png}
    \caption{L'immagine illustra l'approssimazione di un cerchio con poliedri
    convessi.}
    \label{fig:th_07_02}
\end{figure}

Mancando un elemento minimo nell'astratto che sovra-approssima il cerchio,
non è possibile definire $\alpha(P)$ in modo unico come richiesto dalla
definizione di Galois Connection (\Cref{def:galois_connection}).
\end{esempio}

\subsection{Estensione agli stati}

Finora si è parlato di astrazione di valori (es. un intero).
Ma le analisi lavorano su stati ($\Sigma: \text{Var} \to \text{Valori}$).

\begin{proposizione}{Astrazione pointwise}{pointwise_abstraction}

Se esiste una GC tra i valori concreti e i valori astratti, allora
l'applicazione puntuale (pointwise) di questa astrazione agli stati
forma anch'essa una Connessione di Galois.

Questo permette di costruire l'analisi componendo le astrazioni delle singole
variabili.
\end{proposizione}
